% 合成示例简历（保研场景）：人物、机构与数字全部虚构，仅用于演示场景归类与警示规则。
% 与求职版同组（简历多源），按 sources.json 的「简历场景映射」归入保研版场景。
\documentclass{article}
\newcommand{\resumesection}[1]{\section*{#1}}
\newcommand{\resumeentry}[2]{\textbf{#1}\hfill{#2}}
\newcommand{\resumesubentry}[2]{\textit{#1}\hfill{#2}}
\begin{document}

\resumesection{教育背景}
\resumeentry{某大学 · 社会学硕士}{2024 年 9 月 -- 2027 年 6 月（预期）}
\item GPA：3.8/4.0，专业排名 5/30。
\item 核心课程：计量经济学（92）、时间序列分析（92）。

\resumesection{科研经历}
\resumeentry{创业者跟踪调查数据整理}{2025 年 3 月 -- 至今}
\resumesubentry{研究助理}{某大学}
\item 负责问卷数据清洗与校验，覆盖 8,974 个观测。
\item 整理 16--60 岁创业者样本的访谈纪要 40 份。

\resumesection{技能}
\item 工具语言：英语（CET-6：600，IELTS：7.5）。
\item 技术栈：Stata、Python；熟悉面板数据与因果推断方法。

\end{document}